
\documentclass{article}
\usepackage{colortbl}
\usepackage{makecell}
\usepackage{multirow}
\usepackage{supertabular}

\begin{document}

\newcounter{utterance}

\centering \large Interaction Transcript for game `cladder', experiment `full\_v1.5\_default', episode 2646 with refineonly{-}Qwen3.5{-}9B.
\vspace{24pt}

{ \footnotesize  \setcounter{utterance}{1}
\setlength{\tabcolsep}{0pt}
\begin{supertabular}{c@{$\;$}|p{.15\linewidth}@{}p{.15\linewidth}p{.15\linewidth}p{.15\linewidth}p{.15\linewidth}p{.15\linewidth}}
    \# & \multicolumn{2}{c}{Player} && \multicolumn{2}{c}{Game Master} \\
    \hline

    \theutterance \stepcounter{utterance}  
    & & & \multicolumn{4}{p{0.6\linewidth}}{
        \cellcolor[rgb]{0.9,0.9,0.9}{
            \makecell[{{p{\linewidth}}}]{
                \texttt{\tiny{[P1$\langle$GM]}}
                \texttt{You are an expert in causal inference. The following question is not a typical commonsense query, but rather a meticulously designed question created by a professor specializing in causal inference, intended to assess the students' mastery of the course content.} \\
\texttt{Imagine a self{-}contained, hypothetical world with only the following conditions, and without any unmentioned factors or causal relationships: Gender has a direct effect on smoking and lung cancer. Smoking has a direct effect on tar deposit. Tar deposit has a direct effect on lung cancer. Gender is unobserved. For nonsmokers and with no tar deposit, the probability of lung cancer is 79\%. For nonsmokers and with high tar deposit, the probability of lung cancer is 49\%. For smokers and with no tar deposit, the probability of lung cancer is 71\%. For smokers and with high tar deposit, the probability of lung cancer is 40\%. For nonsmokers, the probability of high tar deposit is 79\%. For smokers, the probability of high tar deposit is 40\%. For smokers, would it be more likely to see lung cancer if the person had been a nonsmoker?} \\
\texttt{Start your answer with 'Yes' or 'No', followed by additional reasoning or evidence to support your explanation.} \\
            }
        }
    }
     \\ \\

    \theutterance \stepcounter{utterance}  
    & \multicolumn{4}{p{0.6\linewidth}}{
        \cellcolor[rgb]{0.561,0.737,0.561}{
            \makecell[{{p{\linewidth}}}]{
                \texttt{\tiny{[P1$\rangle$GM]}}
                \texttt{No} \\
\\ 
\texttt{**Reasoning:**} \\
\\ 
\texttt{To determine if a smoker would be more likely to see lung cancer if they had been a nonsmoker, we must compare the observed probability of lung cancer for smokers, \$P(L | S)\$, with the counterfactual probability \$P(L | do(\textbackslash{}neg S), S)\$, which represents the probability of lung cancer for a smoker had they been intervened upon to become a nonsmoker.} \\
\\ 
\texttt{We are given the following conditional probabilities:} \\
\texttt{1.  \$P(L | \textbackslash{}neg S, \textbackslash{}neg T) = 0.79\$} \\
\texttt{2.  \$P(L | \textbackslash{}neg S, T) = 0.49\$} \\
\texttt{3.  \$P(L | S, \textbackslash{}neg T) = 0.71\$} \\
\texttt{4.  \$P(L | S, T) = 0.40\$} \\
\texttt{5.  \$P(T | \textbackslash{}neg S) = 0.79\$ (Probability of high tar for nonsmokers)} \\
\texttt{6.  \$P(T | S) = 0.40\$ (Probability of high tar for smokers)} \\
\\ 
\texttt{**Step 1: Calculate the observed probability of lung cancer for smokers (\$P(L | S)\$).**} \\
\texttt{Using the law of total probability based on the actual distribution of tar deposits among smokers:} \\
\texttt{\$\$P(L | S) = P(L | S, T)P(T | S) + P(L | S, \textbackslash{}neg T)P(\textbackslash{}neg T | S)\$\$} \\
\texttt{\$\$P(L | S) = (0.40 \textbackslash{}times 0.40) + (0.71 \textbackslash{}times (1 {-} 0.40))\$\$} \\
\texttt{\$\$P(L | S) = 0.16 + (0.71 \textbackslash{}times 0.60)\$\$} \\
\texttt{\$\$P(L | S) = 0.16 + 0.426 = 0.586\$\$} \\
\\ 
\texttt{**Step 2: Calculate the counterfactual probability of lung cancer for smokers if they were nonsmokers (\$P(L | do(\textbackslash{}neg S), S)\$).**} \\
\texttt{Under the intervention \$do(\textbackslash{}neg S)\$, the causal mechanism generating tar deposits changes to that of the nonsmoker population. Therefore, the probability of having high tar deposit becomes \$P(T | \textbackslash{}neg S) = 0.79\$. The relationship between tar deposit and lung cancer remains unchanged.} \\
\texttt{\$\$P(L | do(\textbackslash{}neg S), S) = P(L | \textbackslash{}neg S, T)P(T | \textbackslash{}neg S) + P(L | \textbackslash{}neg S, \textbackslash{}neg T)P(\textbackslash{}neg T | \textbackslash{}neg S)\$\$} \\
\texttt{\$\$P(L | do(\textbackslash{}neg S), S) = (0.49 \textbackslash{}times 0.79) + (0.79 \textbackslash{}times (1 {-} 0.79))} \\
            }
        }
    }
    & & \\ \\

    \theutterance \stepcounter{utterance}  
    & & & \multicolumn{2}{p{0.3\linewidth}}{
        \cellcolor[rgb]{0.9,0.9,0.9}{
            \makecell[{{p{\linewidth}}}]{
                \texttt{\tiny{[GM$|$GM]}}
                \texttt{No} \\
            }
        }
    }
    & & \\ \\

    \theutterance \stepcounter{utterance}  
    & & & \multicolumn{2}{p{0.3\linewidth}}{
        \cellcolor[rgb]{0.9,0.9,0.9}{
            \makecell[{{p{\linewidth}}}]{
                \texttt{\tiny{[GM$|$GM]}}
                \texttt{no} \\
            }
        }
    }
    & & \\ \\

    \theutterance \stepcounter{utterance}  
    & & & \multicolumn{2}{p{0.3\linewidth}}{
        \cellcolor[rgb]{0.9,0.9,0.9}{
            \makecell[{{p{\linewidth}}}]{
                \texttt{\tiny{[GM$|$GM]}}
                \texttt{game\_result = WIN} \\
            }
        }
    }
    & & \\ \\

\end{supertabular}
}

\end{document}
